\documentclass[10pt]{article}
% Memo style embedded from /Users/jvi019/src/mm_memos/mmmemo.sty.
\makeatletter
%% mmmemo.sty -- compact technical-memorandum style
%% Inspired by the direct, utilitarian structure of the MIT Haystack
%% Observatory EDGES memo series. No MIT or Haystack branding is reproduced.

\RequirePackage[letterpaper,top=0.58in,bottom=0.62in,left=0.72in,right=0.72in]{geometry}
\RequirePackage[T1]{fontenc}
\RequirePackage{newtxtext}
\RequirePackage{newtxmath}
\RequirePackage{microtype}
\RequirePackage{graphicx}
\RequirePackage{booktabs}
\RequirePackage{siunitx}
\RequirePackage{amsmath}
\RequirePackage{enumitem}
\RequirePackage{needspace}
\RequirePackage{caption}
\RequirePackage{fancyhdr}
\RequirePackage[hidelinks]{hyperref}
\RequirePackage{xcolor}

\definecolor{mmrule}{gray}{0.25}
\setlength{\parindent}{0pt}
\setlength{\parskip}{0.48em plus 0.12em minus 0.08em}
\setlist{nosep,leftmargin=1.35em}
\setlength{\textfloatsep}{10pt plus 2pt minus 2pt}
\setlength{\floatsep}{9pt plus 2pt minus 2pt}
\captionsetup{font=small,labelfont=bf,skip=4pt}
\sisetup{detect-all,per-mode=symbol,range-phrase=--}

\newif\ifscriptprovenance
\scriptprovenancetrue
\newcommand{\showscriptprovenance}{\scriptprovenancetrue}
\newcommand{\hidescriptprovenance}{\scriptprovenancefalse}
\newcommand{\scriptprovenance}[1]{%
  \ifscriptprovenance\par\vspace{1pt}{\footnotesize\itshape Script: \texttt{#1}.}\fi}

\newcommand{\mmseries}{MICROMETEORITE MODELING MEMO SERIES}
\newcommand{\mminstitution}{UIT THE ARCTIC UNIVERSITY OF NORWAY}
\newcommand{\mmlocation}{TROMS\O, NORWAY}
\newcommand{\mmmemonumber}{}
\newcommand{\mmdate}{}
\newcommand{\mmto}{}
\newcommand{\mmfrom}{}
\newcommand{\mmsubject}{}
\newcommand{\mmlogofile}{assets/uit-circular-logo.pdf}
\newcommand{\setmmlogo}[1]{\renewcommand{\mmlogofile}{#1}}

\newcommand{\mmmemo}[5]{%
  \renewcommand{\mmmemonumber}{#1}%
  \renewcommand{\mmdate}{#2}%
  \renewcommand{\mmto}{#3}%
  \renewcommand{\mmfrom}{#4}%
  \renewcommand{\mmsubject}{#5}%
}

\newcommand{\makemmmemo}{%
  \thispagestyle{fancy}%
  \begin{minipage}[t]{0.62\textwidth}
    \begin{minipage}[c]{0.16\textwidth}
      \IfFileExists{\mmlogofile}{\includegraphics[width=0.62in]{\mmlogofile}}{}
    \end{minipage}\hspace{0.02\textwidth}%
    \begin{minipage}[c]{0.80\textwidth}
      \raggedright\bfseries
      \mminstitution\\[-1pt]
      \mmseries\\[-1pt]
      \mmlocation
    \end{minipage}
  \end{minipage}%
  \hfill
  \begin{minipage}[t]{0.34\textwidth}
    \raggedleft\bfseries MM MEMO \#\mmmemonumber\\[2pt]
    \normalfont\mmdate
  \end{minipage}
  \par\vspace{0.42em}\hrule height 0.7pt\vspace{0.48em}
  \begin{tabular}{@{}p{0.72in}p{\dimexpr\textwidth-0.72in-2\tabcolsep\relax}@{}}
    \textbf{To:} & \mmto\\
    \textbf{From:} & \mmfrom\\
    \textbf{Subject:} & \textbf{\mmsubject}
  \end{tabular}
  \par\vspace{0.35em}\hrule height 0.35pt\vspace{0.5em}
}

\pagestyle{fancy}
\fancyhf{}
\fancyhead[L]{\footnotesize MM MEMO \#\mmmemonumber}
\fancyhead[R]{\footnotesize\mmsubject}
\fancyfoot[C]{\footnotesize\thepage}
\renewcommand{\headrulewidth}{0.35pt}
\renewcommand{\footrulewidth}{0pt}
\setlength{\headheight}{13.5pt}

\newcommand{\mmsection}[1]{%
  \Needspace{4\baselineskip}\par\vspace{0.45em}{\bfseries #1}\par\vspace{-0.12em}%
}

\newenvironment{mmconclusion}
  {\par\medskip\hrule\smallskip\textbf{Conclusion.}\ }
  {\par\smallskip\hrule\medskip}

\makeatother

\IfFileExists{memos/assets/uit-circular-logo.pdf}{\setmmlogo{memos/assets/uit-circular-logo.pdf}}{\setmmlogo{assets/uit-circular-logo.pdf}}
\mmmemo{005}{October 11, 2026}{Micrometeorite modeling group}{Juha-Pekka Vierinen}{Zenith-angle normalization of ablation peak heights}
% Toggle implementation provenance with \hidescriptprovenance.
\begin{document}
\makemmmemo

\textbf{Proposal.} Normalize the altitude of peak mass loss to a common
trajectory zenith angle of $45^\circ$. An equal atmospheric-column correction
is a useful first approximation. A correction table from the full thermal
ablation model performs better in the numerical experiment below. Only
ablation is calculated; interpreting initial radar detection heights remains
a later observational test.

\mmsection{Definition and first-order correction}
Let $\chi$ be the \emph{meteor trajectory} angle from local vertical, not the
radar beam zenith angle. Write $P=-dm/dt>0$ and define
$h_p(\chi)$ as the height where $P$ reaches its maximum. For a plane-parallel
atmosphere with mass density $\rho(h)$, define the vertical and traversed
slant mass columns, in \si{kg.m^{-2}}, as
\begin{equation}
 X(h)=\int_h^\infty\rho(z)\,dz,\qquad
 \Sigma(h,\chi)=\frac{X(h)}{\cos\chi}.
\end{equation}
If the peak occurs at approximately the same traversed column for all angles,
map an observed or modelled peak to
\begin{equation}\label{eq:column}
 \boxed{h_{p,45}=X^{-1}\!\left[X(h_p)\frac{\cos45^\circ}{\cos\chi}\right].}
\end{equation}
For $\rho(h)\propto\exp(-h/H)$, this becomes
\begin{equation}\label{eq:exp}
 \boxed{h_{p,45}=h_p+H\ln\!\left(\frac{\cos\chi}{\cos45^\circ}\right).}
\end{equation}
An illustrative $H=\SI{6}{km}$ shifts vertical-entry peaks upward by
\SI{2.08}{km}, leaves $45^\circ$ peaks unchanged, and shifts $60^\circ$
peaks downward by \SI{2.08}{km}. For MSIS, use the numerical column in
Eq.~\eqref{eq:column} rather than one constant scale height.

The approximation has a physical basis. In an energy-limited model,
$P=\Lambda A\rho v^3/(2L)$ and $d\Sigma/dt=\rho v$. Neglecting gravity,
$dm/d\Sigma$ and $dv/d\Sigma$ have no explicit $\chi$ dependence. In an
exponential atmosphere, $\rho=\Sigma\cos\chi/H$, so
$P(\Sigma,\chi)=\cos\chi\,F(\Sigma)$ for otherwise identical grains.
The maximizing column is independent of angle, giving
Eq.~\eqref{eq:exp}. Its ideal peak-rate scaling is also
$P_{p,45}\simeq P_p\cos45^\circ/\cos\chi$.
The full thermal model includes finite heating time, radiative cooling,
temperature-dependent evaporation and changing scale height; therefore the
column equivalence need not be exact.

\newpage
\mmsection{Check using the native Rust ablation model}
We compared 100, 300, 500, 700, 900 and 1100~\si{\micro m} diameters;
20, 40 and \SI{70}{km.s^{-1}} initial speeds; and zenith angles
$0,15,30,45,60,70^\circ$: 108 trajectories. All use the same stony thermal
parameters, \SI{3300}{kg.m^{-3}} bulk density, and MSIS~2.1 atmosphere at
$69.3^\circ$~N, $16.04^\circ$~E on June 28, 2018, 12:45:33 UTC.
Peaks are locally interpolated from samples separated by at most
\SI{25}{m} in height. Each case is compared with its own actual $45^\circ$
model peak, holding diameter and initial speed fixed.

\begin{center}
\begin{tabular}{lrr}
\toprule
Correction & RMS height error (km) & Maximum error (km)\\
\midrule
None & 1.685 & 4.578\\
Equal MSIS mass column & 0.484 & 2.148\\
Model table, held-out diameter & 0.159 & 0.525\\
\bottomrule
\end{tabular}
\end{center}
\scriptprovenance{examples/normalize\_ablation\_zenith.py; src/bin/meteor-ablation.rs}

Thus the column correction reduces pooled RMS differences by a factor
of about 3.5, but does not remove all angle dependence. The worst corrected
case is a \SI{100}{\micro m} grain at \SI{70}{km.s^{-1}} and
$70^\circ$ entry. Halving its height step to \SI{12.5}{m}, and that of its
$45^\circ$ reference, shifts either peak by less than \SI{0.8}{m}; its
\SI{2.15}{km} residual is not explained by peak sampling.

\mmsection{Preferred correction for a height--speed histogram}
For each fixed $v_0,d,\chi$, compute $h_p$ from ablation alone and tabulate
\begin{equation}
 \Delta h(v_0,d,\chi)=h_p(v_0,d,45^\circ)-h_p(v_0,d,\chi).
\end{equation}
Apply $h_{p,45}=h_p+\Delta h$. If diameter is unknown, use a stated
population weighting over diameter, or a provisional median correction
$\widehat{\Delta h}(v_0,\chi)$. The last table row predicts each diameter
using the median correction from the other five diameters. It tests size
transfer within these model cases; interpolation in speed, different materials
and real detection heights have not been validated. Equal weighting of these
six nodes is not an inferred meteoroid size distribution.

For neutral-density studies, retain the original heights and use one specified
reference atmosphere consistently across the samples being compared. In the
exponential limit, a multiplicative density change $f$ shifts the normalized
peak by $\delta h=H\ln f$; this geometric normalization preserves that signal.
Measure the width of the corrected height distribution within narrow speed
bins and check for remaining zenith-angle trends. Angular uncertainty alone
contributes approximately $H\tan\chi\,\sigma_\chi$ to height uncertainty
($\sigma_\chi$ in radians). Near-horizontal trajectories require curved-path
columns; the plane-parallel test here stops at $70^\circ$.

\textbf{Recommendation.} Start with Eq.~\eqref{eq:column}, then use the
ablation-derived correction table if its residual angle dependence is smaller.
The result supports a tighter \emph{model ablation-peak} distribution;
reduction of the observed initial-detection-height width is still to be tested.

\vspace{2pt}
{\footnotesize Atmospheric model: Emmert et al. (2021),
\href{https://doi.org/10.1029/2020EA001321}{NRLMSIS 2.0} (MSIS 2.1 used here).
Observational motivation: Huyghebaert et al. (2026),
\href{https://doi.org/10.5194/amt-19-4277-2026}{MAARSY neutral-density variations}.}
\end{document}
